%%^^A%%  fontspec-doc-opentype.tex -- part of FONTSPEC <latex3.github.io/fontspec>


\documentclass[a4paper]{l3doc}
\usepackage{fontspec-doc-style}
\showexamplesfalse
\begin{document}

\part{OpenType}
\label{sec:opentype-features}

\section{引言}
\label{sec:opentype-features-intro}

OpenType 字体（以及 AAT、Graphite 等其他“智能”字体技术）可以通过多种方式
改变文字外观。这些变化被称为字体特性。当用户把某一特性（例如小型大写）应用于
一段文字时，字体内部的代码会做出相应替换，以小型大写取代小写字母。
但使用这类特性并不会影响底层文本。在我们的小型大写示例中，小写字母仍存储在
文档中；只是外观被 OpenType 特性改变了。这使得文本可以毫无障碍地被检索和复制。
若用户换用另一种不支持小型大写的字体，则会显示“普通”小写字母。

有些 OpenType 特性是支持特定文字所必需的，这些特性通常会自动应用。
例如印度系文字往往要求在用户输入字符后重新塑形、重新排序，
以便按读者预期的传统方式显示。另一些特性可用于支持特定语言。
面向中世纪研究者的 Junicode 字体默认使用古英语形态的字母 thorn，
而现代冰岛语中的 thorn 形态更圆润。若用户把某段文字标记为冰岛语，
Junicode 会通过一个本地化字母形态的 OpenType 特性自动切换为冰岛语形态。

还有一大类 OpenType 特性，旨在为众多语言与文字体系提供高质量排版支持。
前面各节已展示过部分字体特性示例；\pkg{fontspec} 支持的完整 OpenType
字体特性集合在下文 \ref{sec:ot-feat} 中说明。

OpenType 规范为每个特性规定了四字母代码（例如小型大写为 \texttt{smcp}）。
下文在列出各特性的 \pkg{fontspec} 名称时一并给出四字母代码，
供已熟悉 OpenType 的读者参考；若这些代码对你没有意义，可忽略它们。


\subsection{如何选择字体特性}

字体特性通过一系列 \meta{feature}=\meta{option} 选择来指定。特性（通常）按逻辑
分组；例如，所有与连字相关的字体特性都通过书写 \verb|Ligatures={...}| 并传入
相应参数访问，参数可以是 \texttt{TeX}、\texttt{Rare} 等，如下文
\ref{sec:ot-feat-liga} 所示。

任何接受非数值输入的特性都可同时给出多个选项，但这样做并非总能生效。
有些选项会以通常显而易见的方式覆盖其他选项；\Verb|Numbers={OldStyle,Lining}|
意义不大，因为这两个选项相互排斥，\XeTeX\ 只会使用最后指定的选项
（本例中以 \opt{Lining} 取代 \opt{OldStyle}）。

若请求了字体所不具备的特性或选项，会在控制台输出中给出警告。
如 \vref{sec:quiet-warnings} 所述，这些警告可通过选择宏包选项
\texttt{[quiet]} 抑制。

\subsection{如何获知我的字体支持哪些字体特性？}

尽管我一直希望在 \pkg{fontspec} 中提供一个显示字体所含 OpenType 特性的功能，
但它始终不在我的优先事项前列。原因之一是已有文档 |opentype-info.tex|，
可在 \textsc{ctan} 上获取，或在终端窗口中输入 |kpsewhich opentype-info.tex| 找到。
复制此文件并放到方便的位置，然后在常用的 \TeX\ 编辑器中打开，
把字体名改为你想查询的字体；用 plain \XeTeX\ 运行后，输出的 \textsc{pdf}
大致如下：

\begin{framed}

%%% From OpenType-info.tex %%%
\DeleteShortVerb\"
\def\myfontname{[Asana-Math.otf]}
\font\testfont="\myfontname/OT" at 10pt\relax

\rightskip=0pt plus 1fil

\font\titlefont=ec-lmssbx10 at 12pt
\font\tenrm=ec-lmss10 at 9pt \tenrm
\font\tentt=ec-lmtt10 at 9pt

\def\fourcharcode#1{\begingroup
 \count0=#1\count1=\count0
 \ifnum\count0=0
  $\langle$default$\rangle$%
 \else
  \tentt
  '%
  \divide\count0 by "1000000
  \char\count0
  \multiply\count0 by "1000000
  \advance\count1 by -\count0
  \count0=\count1
  \divide\count0 by "10000
  \char\count0
  \multiply\count0 by "10000
  \advance\count1 by -\count0
  \count0=\count1
  \divide\count0 by "100
  \char\count0
  \multiply\count0 by "100
  \advance\count1 by -\count0
  \ifnum\count1=32 \ \else \char\count1 \fi
  '%
 \fi
 \endgroup
}

\newcount\scriptcount
\newcount\scriptindex
\newcount\scripttag
\newcount\langcount
\newcount\langindex
\newcount\langtag
\newcount\featurecount
\newcount\featureindex
\newcount\featuretag

\leftline{\titlefont OpenType Layout features found in `\myfontname'}
\nobreak

\scriptcount=\XeTeXOTcountscripts\testfont
\ifnum\scriptcount=0 \noindent None\par\fi

\loop
 \ifnum\scriptindex<\scriptcount
  \smallskip
  \scripttag=\XeTeXOTscripttag\testfont\scriptindex
  \noindent script = \fourcharcode{\scripttag}\endgraf\nobreak
  \langcount=\XeTeXOTcountlanguages\testfont\scripttag
  \advance\langcount by 1 % one extra to get the default language system
  {\loop
    \langtag=\XeTeXOTlanguagetag\testfont\scripttag\langindex
    \indent language = \fourcharcode{\langtag}\endgraf\nobreak
    \featurecount=\XeTeXOTcountfeatures\testfont\scripttag\langtag
    {\indent\indent \hangindent=3\parindent \hangafter=1 features = \loop
      \featuretag=\XeTeXOTfeaturetag\testfont\scripttag\langtag\featureindex
      \fourcharcode{\featuretag}
      \advance\featureindex by 1
      \ifnum\featureindex<\featurecount \repeat\endgraf}
    \advance\langindex by 1
  \ifnum\langindex<\langcount \medskip \repeat}
  \advance\scriptindex by 1
\repeat
\end{framed}

\noindent 我特意在上文选了一个设计上所含字体特性很少的字体；
Latin Modern Roman 这类“常规”文字字体包含的特性要多得多，
我不想让本文档过于杂乱。找到字体所含的脚本、语言与特性后，
还需把 OpenType 标签与 \pkg{fontspec} 使用的“逻辑”名称对照核实。

\paragraph{otfinfo}
另一种更简单的方式是使用随 \TeX{}Live 分发的命令行工具 |otfinfo|。
只需在终端窗口中输入，例如：
\begin{Verbatim}
  otfinfo -f `kpsewhich lmromandunh10-oblique.otf`
\end{Verbatim}
结果如下：
\begin{Verbatim}[frame=single]
aalt	Access All Alternates
cpsp	Capital Spacing
dlig	Discretionary Ligatures
frac	Fractions
kern	Kerning
liga	Standard Ligatures
lnum	Lining Figures
onum	Oldstyle Figures
pnum	Proportional Figures
size	Optical Size
tnum	Tabular Figures
zero	Slashed Zero
\end{Verbatim}


\section{OpenType 脚本与语言}\label{sec:ot}

包含多种文字与语言字形的字体，可能针对所支持的不同字符集和语言提供不同的
字体特性，而同一字体特性也可能随所选脚本或语言表现不同。使用多语言字体时，
重要的是选定它用于哪种语言，更重要的是选定使用哪种脚本。

“脚本”指所用的字母体系；例如英语和法语都使用拉丁脚本。类似地，
阿拉伯脚本既可书写阿拉伯语，也可书写波斯语。

\feat{Script} 与 \feat{Language} 特性用于指定这些信息。可选选项分别列于
\vref{tab:ot-scpt} 与 \vref{tab:ot-lang}。若请求了当前字体不支持的脚本或语言，
会在控制台输出中打印警告。如需创建新的 \feat{Script} 或 \feat{Language}
选项，方法见 \vref{sec:newscriptlang}。

由于这些字体特性会改变字体可选哪些特性，\pkg{fontspec} 会在所有其他特性之前
自动选定 \feat{Script} 与 \feat{Language} 设置；若使用 \XeTeX，还会按
\vref{sec:renderer-xetex} 所述专门为该字体选用 \opt{OpenType} 渲染器。

OpenType 字体可根据所选脚本与语言提供不同的字体特性。此外，这些设置还可以
建立自身的字体行为与字形选择（一个例子是保加利亚语、塞尔维亚语与俄语所用
字母表中部分字母的风格差异）。若出于某种原因需要，\pkg{fontspec} 特性
\feat{LocalForms} \texttt{=} \opt{Off} 可禁用其中部分替换。需要注意，
\feat{LocalForms} \texttt{=} \opt{On} 并非 \pkg{fontspec} 的默认设置，
而是 \XeTeX\ 与 \LuaTeX/\pkg{otfload} 底层字体成形引擎的默认设置。


\subsection{\feat{Script} 与 \feat{Language} 示例}

在 \exref{script-lang} 所示示例中，使用 Code2000 字体\note{\url{http://www.code2000.net/}}
为各种字母体系排版不同输入文本，分别对比应用与不应用 OpenType 脚本的效果。
只有第二种情形下文字才正确呈现；可以看到许多附加符号间距不正确、
缺少上下文连字与重排的例子。感谢 \name{Jonathan Kew}、\name{Yves Codet}
与 \name{Gildas Hamel} 对这些示例的贡献。

\begin{Xexample}[firstline=14,lastline=23]{script-lang}{各种脚本与语言示例。}
\def\testfeature#1#2{%^^A
  \fontspec{\examplefont}#2 & \fontspec[#1]{\examplefont}#2\\[1ex]}
\def \arabictext{العربي}
\def \devanagaritext{हिन्दी}
\def \bengalitext{লেখ}
\def \gujaratitext{મર્યાદા-સૂચક નિવેદન}
\def \malayalamtext{നമ്മുടെ പാരബര്യ}
\def \gurmukhitext{ਆਦਿ ਸਚੁ ਜੁਗਾਦਿ ਸਚੁ}
\def \tamiltext{தமிழ் தேடி}
\def \hebrewtext{רִדְתָּֽהּ}
\def \vietnamesetext{cấp số mỗi}
\begin{tabular}{r@{\quad}l}
  \gdef\examplefont{FreeSerif.otf}
  \testfeature{Script=Arabic}{\arabictext}
  \testfeature{Script=Devanagari}{\devanagaritext}
  \testfeature{Script=Bengali}{\bengalitext}
  \testfeature{Script=Gujarati}{\gujaratitext}
  \testfeature{Script=Malayalam}{\malayalamtext}
  \testfeature{Script=Gurmukhi}{\gurmukhitext}
  \testfeature{Script=Tamil}{\tamiltext}
  \testfeature{Script=Hebrew}{\hebrewtext}
  \def\examplefont{Doulos_SIL_Regular.ttf}
  \testfeature{Language=Vietnamese}{\vietnamesetext}
\end{tabular}
\end{Xexample}



\begin{table}[!hbp]
  \caption{为 OpenType 字体定义的 \opt{Script}。别名以红色¶（{\sffamily\textcolor{red}{\P}}）标记，并列于相邻位置。}
  \label{tab:ot-scpt}
\def\dup{\makebox[0pt][r]{\textcolor{red}{\P}}}
\setlength\columnseprule{0pt}
  \hrule
  \begin{multicols}{4}\setlength\parindent{0pt}
    \sffamily\scriptsize
    Adlam \par
    Ahom \par
    Anatolian Hieroglyphs \par
    Arabic \par
    Armenian \par
    Avestan \par
    Balinese \par
    Bamum \par
    Bassa Vah \par
    Batak \par
    Bengali \par
    Bhaiksuki \par
    Bopomofo \par
    Brahmi \par
    Braille \par
    Buginese \par
    Buhid \par
    Byzantine Music \par
    Canadian Syllabics \par
    Carian \par
    Caucasian Albanian \par
    Chakma \par
    Cham \par
    Cherokee \par
    \dup CJK \par
    \dup CJK Ideographic \par
    Coptic \par
    Cypriot Syllabary \par
    Cyrillic \par
    Default \par
    Deseret \par
    Devanagari \par
    Duployan \par
    Egyptian Hieroglyphs \par
    Elbasan \par
    Ethiopic \par
    Georgian \par
    Glagolitic \par
    Gothic \par
    Grantha \par
    Greek \par
    Gujarati \par
    Gurmukhi \par
    Hangul Jamo \par
    Hangul \par
    Hanunoo \par
    Hatran \par
    Hebrew \par
    \dup Hiragana and Katakana \par
    \dup Kana \par
    Imperial Aramaic \par
    Inscriptional Pahlavi \par
    Inscriptional Parthian \par
    Javanese \par
    Kaithi \par
    Kannada \par
    Kayah Li \par
    Kharosthi \par
    Khmer \par
    Khojki \par
    Khudawadi \par
    Lao \par
    Latin \par
    Lepcha \par
    Limbu \par
    Linear A \par
    Linear B \par
    Lisu \par
    Lycian \par
    Lydian \par
    Mahajani \par
    Malayalam \par
    Mandaic \par
    Manichaean \par
    Marchen \par
    \dup Math \par
    \dup Maths \par
    Meitei Mayek \par
    Mende Kikakui \par
    Meroitic Cursive \par
    Meroitic Hieroglyphs \par
    Miao \par
    Modi \par
    Mongolian \par
    Mro \par
    Multani \par
    Musical Symbols \par
    Myanmar \par
    \dup N'Ko \par
    \dup N'ko \par
    Nabataean \par
    Newa \par
    Ogham \par
    Ol Chiki \par
    Old Italic \par
    Old Hungarian \par
    Old North Arabian \par
    Old Permic \par
    Old Persian Cuneiform \par
    Old South Arabian \par
    Old Turkic \par
    \dup Oriya \par
    \dup Odia \par
    Osage \par
    Osmanya \par
    Pahawh Hmong \par
    Palmyrene \par
    Pau Cin Hau \par
    Phags-pa \par
    Phoenician \par
    Psalter Pahlavi \par
    Rejang \par
    Runic \par
    Samaritan \par
    Saurashtra \par
    Sharada \par
    Shavian \par
    Siddham \par
    Sign Writing \par
    Sinhala \par
    Sora Sompeng \par
    Sumero-Akkadian Cuneiform \par
    Sundanese \par
    Syloti Nagri \par
    Syriac \par
    Tagalog \par
    Tagbanwa \par
    Tai Le \par
    Tai Lu \par
    Tai Tham \par
    Tai Viet \par
    Takri \par
    Tamil \par
    Tangut \par
    Telugu \par
    Thaana \par
    Thai \par
    Tibetan \par
    Tifinagh \par
    Tirhuta \par
    Ugaritic Cuneiform \par
    Vai \par
    Warang Citi \par
    Yi
  \end{multicols}
  \hrule
\end{table}

\begin{table}[p]
  \vspace*{-3cm}
  \hspace{-3cm}
  \def\dup{\makebox[0pt][r]{\textcolor{red}{\P}}}
  \begin{minipage}{\linewidth+4cm}
  \caption{为 OpenType 字体定义的 \opt{Language}。别名以红色¶（{\sffamily\textcolor{red}{\P}}）标记，并列于相邻位置。}
  \label{tab:ot-lang}
  \setlength\columnseprule{0pt}
  \hrule
  \begin{multicols}{6}
    \everypar{\setlength\parindent{0pt}\setlength\hangindent{2em}}
    \sffamily\footnotesize\raggedright
    Abaza \par
    Abkhazian \par
    Adyghe \par
    Afrikaans \par
    Afar \par
    Agaw \par
    Altai \par
    Amharic \par
    Arabic \par
    Aari \par
    Arakanese \par
    Assamese \par
    Athapaskan \par
    Avar \par
    Awadhi \par
    Aymara \par
    Azeri \par
    Badaga \par
    Baghelkhandi \par
    Balkar \par
    Baule \par
    Berber \par
    Bench \par
    Bible Cree \par
    Belarussian \par
    Bemba \par
    Bengali \par
    Bulgarian \par
    Bhili \par
    Bhojpuri \par
    Bikol \par
    Bilen \par
    Blackfoot \par
    Balochi \par
    Balante \par
    Balti \par
    Bambara \par
    Bamileke \par
    Breton \par
    Brahui \par
    Braj Bhasha \par
    Burmese \par
    Bashkir \par
    Beti \par
    Catalan \par
    Cebuano \par
    Chechen \par
    Chaha Gurage \par
    Chattisgarhi \par
    Chichewa \par
    Chukchi \par
    Chipewyan \par
    Cherokee \par
    Chuvash \par
    Comorian \par
    Coptic \par
    Cree \par
    Carrier \par
    Crimean Tatar \par
    Church Slavonic \par
    Czech \par
    Danish \par
    Dargwa \par
    Woods Cree \par
    German \par
    Default \par
    Dogri \par
    Divehi \par
    Djerma \par
    Dangme \par
    Dinka \par
    Dungan \par
    Dzongkha \par
    Ebira \par
    Eastern Cree \par
    Edo \par
    Efik \par
    Greek \par
    English \par
    Erzya \par
    Spanish \par
    Estonian \par
    Basque \par
    Evenki \par
    Even \par
    Ewe \par
    French Antillean \par
    \dup Farsi \par
    \dup Parsi \par
    \dup Persian \par
    Finnish \par
    Fijian \par
    Flemish \par
    Forest Nenets \par
    Fon \par
    Faroese \par
    French \par
    Frisian \par
    Friulian \par
    Futa \par
    Fulani \par
    Ga \par
    Gaelic \par
    Gagauz \par
    Galician \par
    Garshuni \par
    Garhwali \par
    Ge'ez \par
    Gilyak \par
    Gumuz \par
    Gondi \par
    Greenlandic \par
    Garo \par
    Guarani \par
    Gujarati \par
    Haitian \par
    Halam \par
    Harauti \par
    Hausa \par
    Hawaiin \par
    Hammer-Banna \par
    Hiligaynon \par
    Hindi \par
    High Mari \par
    Hindko \par
    Ho \par
    Harari \par
    Croatian \par
    Hungarian \par
    Armenian \par
    Igbo \par
    Ijo \par
    Ilokano \par
    Indonesian \par
    Ingush \par
    Inuktitut \par
    Irish \par
    Irish Traditional \par
    Icelandic \par
    Inari Sami \par
    Italian \par
    Hebrew \par
    Javanese \par
    Yiddish \par
    Japanese \par
    Judezmo \par
    Jula \par
    Kabardian \par
    Kachchi \par
    Kalenjin \par
    Kannada \par
    Karachay \par
    Georgian \par
    Kazakh \par
    Kebena \par
    Khutsuri Georgian \par
    Khakass \par
    Khanty-Kazim \par
    Khmer \par
    Khanty-Shurishkar \par
    Khanty-Vakhi \par
    Khowar \par
    Kikuyu \par
    Kirghiz \par
    Kisii \par
    Kokni \par
    Kalmyk \par
    Kamba \par
    Kumaoni \par
    Komo \par
    Komso \par
    Kanuri \par
    Kodagu \par
    Korean Old Hangul \par
    Konkani \par
    Kikongo \par
    Komi-Permyak \par
    Korean \par
    Komi-Zyrian \par
    Kpelle \par
    Krio \par
    Karakalpak \par
    Karelian \par
    Karaim \par
    Karen \par
    Koorete \par
    Kashmiri \par
    Khasi \par
    Kildin Sami \par
    Kui \par
    Kulvi \par
    Kumyk \par
    Kurdish \par
    Kurukh \par
    Kuy \par
    Koryak \par
    Ladin \par
    Lahuli \par
    Lak \par
    Lambani \par
    Lao \par
    Latin \par
    Laz \par
    L-Cree \par
    Ladakhi \par
    Lezgi \par
    Lingala \par
    Low Mari \par
    Limbu \par
    Lomwe \par
    Lower Sorbian \par
    Lule Sami \par
    Lithuanian \par
    Luba \par
    Luganda \par
    Luhya \par
    Luo \par
    Latvian \par
    Majang \par
    Makua \par
    Malayalam Traditional \par
    Mansi \par
    Marathi \par
    Marwari \par
    Mbundu \par
    Manchu \par
    Moose Cree \par
    Mende \par
    Me'en \par
    Mizo \par
    Macedonian \par
    Male \par
    Malagasy \par
    Malinke \par
    Malayalam Reformed \par
    Malay \par
    Mandinka \par
    Mongolian \par
    Manipuri \par
    Maninka \par
    Manx Gaelic \par
    Moksha \par
    Moldavian \par
    Mon \par
    Moroccan \par
    Maori \par
    Maithili \par
    Maltese \par
    Mundari \par
    Naga-Assamese \par
    Nanai \par
    Naskapi \par
    N-Cree \par
    Ndebele \par
    Ndonga \par
    Nepali \par
    Newari \par
    Nagari \par
    Norway House Cree \par
    Nisi \par
    Niuean \par
    Nkole \par
    N'ko \par
    Dutch \par
    Nogai \par
    Norwegian \par
    Northern Sami \par
    Northern Tai \par
    Esperanto \par
    Nynorsk \par
    Oji-Cree \par
    Ojibway \par
    Oriya \par
    Oromo \par
    Ossetian \par
    Palestinian Aramaic \par
    Pali \par
    Punjabi \par
    Palpa \par
    Pashto \par
    Polytonic Greek \par
    Pilipino \par
    Palaung \par
    Polish \par
    Provencal \par
    Portuguese \par
    Chin \par
    Rajasthani \par
    R-Cree \par
    Russian Buriat \par
    Riang \par
    Rhaeto-Romanic \par
    Romanian \par
    Romany \par
    Rusyn \par
    Ruanda \par
    Russian \par
    Sadri \par
    Sanskrit \par
    Santali \par
    Sayisi \par
    Sekota \par
    Selkup \par
    Sango \par
    Shan \par
    Sibe \par
    Sidamo \par
    Silte Gurage \par
    Skolt Sami \par
    Slovak \par
    Slavey \par
    Slovenian \par
    Somali \par
    Samoan \par
    Sena \par
    Sindhi \par
    Sinhalese \par
    Soninke \par
    Sodo Gurage \par
    Sotho \par
    Albanian \par
    Serbian \par
    Saraiki \par
    Serer \par
    South Slavey \par
    Southern Sami \par
    Suri \par
    Svan \par
    Swedish \par
    Swadaya Aramaic \par
    Swahili \par
    Swazi \par
    Sutu \par
    Syriac \par
    Tabasaran \par
    Tajiki \par
    Tamil \par
    Tatar \par
    TH-Cree \par
    Telugu \par
    Tongan \par
    Tigre \par
    Tigrinya \par
    Thai \par
    Tahitian \par
    Tibetan \par
    Turkmen \par
    Temne \par
    Tswana \par
    Tundra Nenets \par
    Tonga \par
    Todo \par
    Turkish \par
    Tsonga \par
    Turoyo Aramaic \par
    Tulu \par
    Tuvin \par
    Twi \par
    Udmurt \par
    Ukrainian \par
    Urdu \par
    Upper Sorbian \par
    Uyghur \par
    Uzbek \par
    Venda \par
    Vietnamese \par
    Wa \par
    Wagdi \par
    West-Cree \par
    Welsh \par
    Wolof \par
    Tai Lue \par
    Xhosa \par
    Yakut \par
    Yoruba \par
    Y-Cree \par
    Yi Classic \par
    Yi Modern \par
    Chinese Hong Kong \par
    Chinese Phonetic \par
    Chinese Simplified \par
    Chinese Traditional \par
    Zande \par
    Zulu
  \end{multicols}
  \hspace{4pt}
  \hrule
 \end{minipage}
\end{table}


\section{OpenType 字体特性}
\label{sec:ot-feat}

OpenType 字体特性数量有限，\pkg{fontspec} 为其中约一半提供了接口。
后续各节将给出完整说明，包括如何启用与禁用单个特性，以及各特性如何相互配合。

\vref{tab:all-ot} 提供了一份简要参考；请注意该清单并不完整——只列出了“启用”
键，在为方便起见提供了替代接口时，也只列出第一个。
（例如 |Numbers=OldStyle| 与 |Numbers=Lowercase| 等价。）

为完整起见，\vref{tab:none-ot} 列出了\emph{未}提供 \pkg{fontspec} 接口的
全部 OpenType 特性。遗漏的特性一部分是有意为之，一部分则是疏忽；例如 |aalt|
特性在 \TeX\ 中基本无用，因为它是为通过 \textsc{gui} 界面选择某字形的
“全部替代字形”而设计的。另一些（例如光学边距）则只是由于缺少可用于测试的
字体而尚未考虑。欢迎就如何、在何处把这些缺失特性加入本宏包提出建议。

\ExplSyntaxOn
\def\allOTfeat{
  \prop_map_inline:Nn \g__fontspec_all_opentype_feature_names_prop
    { \opentypefeature{##1}{##2} }
}
\newcommand\opentypefeature[2]{
  \prop_get:NnNT \g__fontspec_OT_features_prop {#1} \tmpa
      {
        \raggedright
        \hangindent=5.2cm
        \makebox[1cm][l]{\textsc{#1}}
        \makebox[4.2cm][l]{
          \int_compare:nT { \tl_count:N \tmpa > 25 } {\ttcondensed} {\ttfamily}
          \tmpa
        }
        \textit{#2}
        \par
        \vspace{2pt}
     }
}
\ExplSyntaxOff

\begin{table}
\caption{\textsf{fontspec} 中 OpenType 特性一览，按特性标签字母序排列。}
\label{tab:all-ot}
\centerline{%
\begin{minipage}{18cm}
\small
\hrule\smallskip
\begin{multicols}{2}
\parindent =0pt
\allOTfeat
\end{multicols}
\vspace*{-\smallskipamount}
\hrule
\end{minipage}}
\end{table}

\ExplSyntaxOn
\renewcommand\opentypefeature[2]{
  \prop_get:NnNF \g__fontspec_OT_features_prop {#1} \tmpa
      {
        \raggedright
        \hangindent=0.9cm
        \makebox[0.9cm][l]{\textsc{#1}}%
        \textit{#2}
        \par
     }
}
\ExplSyntaxOff

\begin{table}
\caption{\emph{不受支持}的 OpenType 特性清单。}
\label{tab:none-ot}
\bigskip
\centerline{%
\begin{minipage}{15cm}
\hrule\smallskip
\begin{multicols}{3}
\parindent =0pt
\allOTfeat
\end{multicols}
\vspace*{-\smallskipamount}
\hrule
\end{minipage}}
\end{table}

\subsection{基于标签的特性}


\subsubsection{替代字形 --- \texttt{salt}}

\feat{Alternate} 特性（别名 \feat{StylisticAlternates}）用于在字体存在变体时
访问替代字体字形，如 \exref{salt} 所示。它采用从 0 开始的数值选择，
具体数值因字体而异。注意 \texttt{Style=Alternate} 等价于
\texttt{Alternate=0}，用于访问默认情形。

\begin{Xexample}[firstline=2]{salt}{\feat{Alternate} 特性。}
  \huge
  \fontspec{LinLibertine_R.otf}
  \textsc{a} \& h \\
  \addfontfeature{Alternate=0}
  \textsc{a} \& h
\end{Xexample}

注意编号从 0 开始。在 \LuaTeX\ 引擎下，|Alternate=Random| 随机选择一个替代字形。
如有需要，为替代字形命名的方式见 \vref{sec:newfeatures}。


\subsubsection{字符变体 --- \texttt{cvNN}}

“字符变体”以数值方式选择，用于调整特定字体中（通常）单个字符的输出。
它们对应 OpenType 特性 |cv01| 到 |cv99|。对于每个可变化的字符，
都可以在该字形的若干可能选项中选择。例如，在下面这个假设示例中，
为字形“4”和“5”选择了变体，末尾的 |:|\meta{n} 对应要选择哪一种变体。
\begin{Verbatim}
  \fontspec{CV Font}[CharacterVariant={4,5:2}] \& violet
\end{Verbatim}
编号完全取决于具体字体；例如字形“5”可能对应字符“v”。字符变体在设计上
互不冲突，因此可以按字符逐个启用（这与风格替代字形不同）。注意编号从 0 开始。


\subsubsection{上下文相关}
此特性指按字形在单词或字符串中的相对位置“随上下文”变化的字形替换；
上下文花体等特性通过 \ref{feat:Contextuals} 所列选项访问。

\begin{features}{Contextuals}
\otf*{Swash}{cswh}
\otf*{Alternate}{calt}
\otf*{WordInitial}{init}
\otf*{WordFinal}{fina}
\otf*{LineFinal}{falt}
\otf*{Inner}{medi}
\cmidrule{2-4}
\otf{ResetAll}{}
\end{features}

OpenType 字体中的历史字形通过 \feat{Style=Historic} 特性访问；
它在 OpenType 中一般\emph{不}属于上下文相关，因此未包含在本特性内。


\subsubsection{附加符号}

指定组合附加符号应如何放置。这些通常会根据脚本设置自动控制。

\begin{features}{Diacritics}
\otf*{MarkToBase}{mark}
\otf*{MarkToMark}{mkmk}
\otf*{AboveBase}{abvm}
\otf*{BelowBase}{blwm}
\cmidrule{2-4}
\otf{ResetAll}{}
\end{features}


\subsubsection{分数 --- \texttt{frac}}

\begin{features}{Fractions}
\otf{On}{+frac}
\otf{Off}{-frac}
\otf{Reset}{}
\cmidrule{2-4}
\otf*{Alternate}{afrc}
\cmidrule{2-4}
\otf{ResetAll}{}
\end{features}

启用“普通”分数的构造：使用字体内部预先组合好的字形，和/或下标与上标字符。
覆盖范围因字体而异；见 \exref{ot-frac}。另有一些字体（主要是亚洲字体）
还提供 \opt{Alternate} 选项。

\begin{Lexample}{ot-frac}{\feat{Fractions} 特性。}
\setsansfont{IBMPlexSans-Regular.otf}[Fractions=On]
\setmonofont{IBMPlexMono-Regular.otf}[Fractions=On]

\sffamily 1/2 47/11 1/1000 \par
\ttfamily 1/2 47/11
\end{Lexample}



\subsubsection{字距调整 --- \texttt{kern}}
\label{sec:kerning}

指定字形间间距如何表现。制作精良的字体包含在不同字符对之间插入不同
间距量的信息。此字距空间会自动插入，但在少数情况下你可能希望将其关闭。

\begin{features}{Kerning}
\otf{On}{+kern}
\otf{Off}{-kern}
\otf{Reset}{}
\cmidrule{2-4}
\otf*{Uppercase}{cpsp}
\cmidrule{2-4}
\otf{ResetAll}{}
\end{features}

如 \vref{sec:letters} 末尾简要提及，\opt{Uppercase} 选项会在大写字母之间
增加少量字距，见 \exref{kernup}，该示例使用 Venturis ADF 字体\note{\url{http://arkandis.tuxfamily.org/adffonts.html}}
（感谢 Clea F. Rees 提供建议）。\opt{Uppercase} 选项与由 \opt{On}/\opt{Off}
选项控制的常规字距调整相互独立。

\begin{Xexample}[firstline=2]{kernup}{为大写字母增加额外字距。（差异通常非常小。）}
  \large
  \fontspec{VenturisADF-Bold.otf}
   UPPERCASE EXAMPLE \\
  \addfontfeature{Kerning=Uppercase}
   UPPERCASE EXAMPLE
\end{Xexample}


\subsubsection{字母} \label{sec:letters}
\opt{Letters} 特性指定当前字体中字母的外观。OpenType 字体可包含下列选项：
\opt{SmallCaps}、\opt{PetiteCaps}、\opt{UppercaseSmallCaps}、
\opt{UppercasePetiteCaps} 与 \opt{Unicase}。此外，在 \LuaTeX\ 中
所有字体都支持 \opt{Uppercase} 与 \opt{Lowercase}。

与早期版本不同，\opt{Uppercase} 与 \opt{Lowercase} 选项会把文字转为大写或
小写，并不要求文字本来就具有相应大小写。\opt{Uppercase} 的旧行为可通过
\verb|Style=Uppercase| 获得。选用 \opt{Uppercase} 选项时，若字体支持，
会自动应用 \verb|Style=Uppercase| 与 \verb|Kerning=Uppercase|。


\begin{features}{Letters}
\otf*{SmallCaps}{smcp}
\otf*{PetiteCaps}{pcap}
\otf*{UppercaseSmallCaps}{c2sc}
\otf*{UppercasePetiteCaps}{c2pc}
\otf*{Unicase}{unic}
\otf*{Uppercase}{}
\otf*{Lowercase}{}
\cmidrule{2-4}
\otf{ResetAll}{}
\end{features}

小形大写（petite caps）比小型大写（small caps）更小。
\opt{SmallCaps} 与 \opt{PetiteCaps} 把小写字母转为较小的大写字母；
而 \opt{Uppercase...} 选项把\emph{大写}字母转为较小的大写字母
（适合应用于本来就全大写的首字母缩写词，如“NASA”）。这一差异见 \exref{caps}。
“Unicase”是大写与小写字母的一种奇特混合。

\begin{Lexample}{caps}{由小写或大写字母生成小型大写。}
  \fontspec{Coelacanth.otf}[Letters=SmallCaps]
   THIS SENTENCE no verb                \\
  \fontspec{Coelacanth.otf}[Letters=UppercaseSmallCaps]
   THIS SENTENCE no verb                \\
  \fontspec{Coelacanth.otf}[Letters=PetiteCaps]
   THIS SENTENCE no verb
\end{Lexample}



\subsubsection{连字}
\label{sec:ot-feat-liga}

\feat{Ligatures} 指出于功能或 \ae sthetic 原因，用专门绘制的字形替换两个独立字符。
可同时选择多个选项，清单见 \ref{feat:Ligatures}。使用 Linux Libertine 字体\note{\url{http://www.linuxlibertine.org/}}
的演示见 \exref{lig}。

注意 \verb|Ligatures=TeX| 所访问的附加特性。它们实际上并非真正的 OpenType
特性，而是 \pkg{luaotfload} 提供的附加功能（即仅 \LuaTeX），用于模拟 \TeX\
对弯引号与标点的 \textsc{ascii} 输入行为。在 \XeTeX\ 中这通过 \feat{Mapping}
特性实现（见 \vref{sec:mapping}）；为保持一致，\verb|Ligatures=TeX| 与
\verb|Mapping=tex-text| 执行相同功能。

\begin{features}{Ligatures}
\otf*{Required}{rlig}
\otf*{Common}{liga}
\otf*{Contextual}{clig}
\otf*{Rare/Discretionary}{dlig}
\otf*{Historic}{hlig}
\otf*{TeX}{tlig}
\cmidrule{2-4}
\otf{ResetAll}{}
\end{features}

\begin{Lexample}[firstline=2]{lig}{\feat{Ligatures} 特性示例。}
   \Huge\centering
   \def\test#1#2{%
     #2 $\to$ {\addfontfeature{#1} #2}\\}
   \fontspec{LinLibertine_R.otf}
   \test{Ligatures=Historic}{strict}
   \test{Ligatures=Rare}{wurtzite}
   \test{Ligatures=CommonOff}{firefly}
\end{Lexample}


\subsubsection{本地化字形 --- \texttt{locl}}

\begin{features}{LocalForms}
\otf{On}{+locl}
\otf{Off}{-locl}
\otf{Reset}{}
\end{features}

此特性用于启用和禁用特定于字体中所选 \feat{Language} 的字形替换等。
该特性存在时默认自动激活，因此一般无需使用 \feat{LocalForms} \texttt{=} \opt{On}。
在某些情形下将其设为 \opt{Off} 可能很重要（不过我一时想不出具体例子）。


\subsubsection{数字}

\feat{Numbers} 特性定义所选字体中数字的外观，接受 \ref{feat:Numbers}
所列选项。

\begin{features}{Numbers}
\otf*{Uppercase}{lnum}
\otf*{Lowercase}{onum}
\otf*{Lining}{lnum}
\otf*{OldStyle}{onum}
\otf*{Proportional}{pnum}
\otf*{Monospaced}{tnum}
\otf*{SlashedZero}{zero}
\otf*{Arabic}{anum}
\cmidrule{2-4}
\otf{ResetAll}{}
\end{features}

同义词 \opt{Uppercase} 与 \opt{Lowercase} 分别等价于
\opt{Lining} 与 \opt{OldStyle}。差异已在 \vref{sec:addfontfeatures} 中展示。
当数字需要纵向对齐时，\opt{Monospaced} 选项对表格材料很有用。

\opt{SlashedZero} 选项把默认的零替换为带斜杠的版本，
以避免与大写“O”混淆，见 \exref{slashzero}。

\begin{Lexample}{slashzero}{\opt{SlashedZero} 选项的效果。}
  \fontspec[Numbers=Lining]{texgyrebonum-regular.otf}
   0123456789
  \fontspec[Numbers=SlashedZero]{texgyrebonum-regular.otf}
   0123456789
\end{Lexample}

\opt{Arabic} 选项（标签 \verb|anum|）根据当前 \opt{Language} 设置
（见 \vref{sec:ot}），把常规数字映射为相应的阿拉伯文或波斯文数字。
此选项基于 \pkg{luaotfload} 宏包的一项 \LuaTeX\ 功能，而非 OpenType 特性。
（因此该特性在 \XeTeX\ 中不可用。）此特性应视为已弃用；虽然目前没有计划
将其从本宏包中移除，但若字体加载器不再支持它，它可能在几乎不加预告的情况下
从 \pkg{fontspec} 中消失。


\subsubsection{装饰图案 --- \texttt{ornm}}

装饰图案通过 \feat{Ornament} 特性（OpenType 特性 |ornm|）选择，
与 \feat{Annotation} 特性一样按数值选取。


\subsubsection{Style}
\label{sec:ot-feat-style}

\begin{features}{Style}
\otf*{Alternate}{salt}
\otf*{Cursive}{curs}
\otf*{Historic}{hist}
\otf*{Italic}{ital}
\otf*{Ruby}{ruby}
\otf*{Swash}{swsh}
\otf*{Titling}{titl}
\otf*{Uppercase}{case}
\otf*{HorizontalKana}{hkna}
\otf*{VerticalKana}{vkna}
\cmidrule{2-4}
\otf{ResetAll}{}
\end{features}

“Ruby”指小型光学字号，在日文排版中用于标注。对于含有多个 |salt|
OpenType 特性的字体，应改用 \pkg{fontspec} 的 \feat{Alternate} 特性。

\Exref{style-alt} 展示了一个涉及字母表中特定字母字形替换的字体特性示例。
这些类别中的其他选项以类似方式运作；具体替换归入哪个特性，很大程度上由字体设计师决定。

 \begin{Xexample}[firstline=2]{style-alt}{\feat{Style} 特性 \opt{Alternate} 选项示例。}
  \Large
  \fontspec{Quattrocento-Regular.ttf}
   M Q W                      \\
  \addfontfeature{Style=Alternate}
   M Q W
\end{Xexample}

\opt{Uppercase} 选项用于为附加符号、破折号等字形选择各种大写形式，
如 \exref{style-uppercase} 所示；注意连字符位置被抬高，以便更好地匹配
周围字母。它（多半）并不会真的把字母映射为大写。
 \note{若需要自动大写字母，请参见 \LaTeX\ 的 \cmd\MakeUppercase\ 命令；
      若使用 \LuaTeX，也可使用 \feat{Letters} 特性。}
此选项过去在 \feat{Letters} 特性下选择，现已移至此处，因为它一般并不实际
影响字母本身。\feat{Kerning} 特性也含有一个 \opt{Uppercase} 选项，
用于在字母之间增加少量间距（见 \vref{sec:kerning}）。

\begin{Lexample}{style-uppercase}{\feat{Style} 特性 \opt{Uppercase} 选项示例。}
  \fontspec{LinLibertine_R.otf}
   UPPER-CASE example \\
  \addfontfeature{Style=Uppercase}
   UPPER-CASE example
\end{Lexample}


另一些特性则能看到更广泛的改变，覆盖整个字母表的风格。
例如在某些日文字体中，|Style=Italic| 或 |Style=Ruby| 等特性分别把所有
拉丁字符改为斜体，或把所有平假名改为较深的光学字形：
\begin{Verbatim}[firstline=2]
  \Large \def\kana{ようこそ ワカヨタレソ}
  \fontspec{Hiragino Mincho Pro}
   Latin \kana        \\
  \addfontfeature{Style={Italic, Ruby}}
   Latin \kana
\end{Verbatim}


\subsubsection{风格集变体 --- \texttt{ssNN}}
\label{sec:ot-ss}

此特性选择某个“风格集”变体，它通常对应一组字符（一般是某个字母表或其子集）
的替代字形风格。该特性以数值指定，对应 OpenType 特性 |ss01|、|ss02| 等。

Junicode 字体\note{\url{http://junicode.sf.net}} 的两个演示见 \exref{ss}
与 \exref{ss2}；感谢 Adam Buchbinder 提供建议。

\begin{Lexample}{ss}{Junicode 字体中使用 \feat{StylisticSet} 特性获得中世纪北欧所用的岛屿字母字形。}
  \fontspec{Junicode}
   Insular forms. \\
  \addfontfeature{StylisticSet=2}
   Insular forms. \\
\end{Lexample}

\begin{Lexample}{ss2}{Junicode 字体中使用 \feat{StylisticSet} 特性获得放大的小写字母（大写字母保持不变）。}
  \fontspec{Junicode}
   ENLARGED Minuscules. \\
  \addfontfeature{StylisticSet=6}
   ENLARGED Minuscules. \\
\end{Lexample}

可同时选择多个风格集，例如写成 |StylisticSet={1,2,3}|。

对于 \AAT\ 字体，|StylisticSet| 特性是 \feat{Variant} 特性的同义词。
如有需要，为风格集命名的方式见 \vref{sec:newfeatures}，这应按字体逐一进行。


\subsubsection{上下位置}

\begin{features}{VerticalPosition}
\otf*{Superior}{sups}
\otf*{Inferior}{subs}
\otf*{Numerator}{numr}
\otf*{Denominator}{dnom}
\otf*{ScientificInferior}{sinf}
\otf*{Ordinal}{ordn}
\cmidrule{2-4}
\otf{ResetAll}{}
\end{features}

\feat{VerticalPosition} 特性用于访问下标（\opt{Inferior}）、上标
（\opt{Superior}）数字与字母（有时还有少量标点）。\opt{Ordinal} 选项只会
提升在某些语言中直接跟在数字之后使用的字符。\opt{ScientificInferior}
特性会把字形移到比 \opt{Inferior} 特性更低的基线以下。这些见 \exref{vertpos}。

\opt{Numerator} 与 \opt{Denominator} 只应用于构造任意分数（见下一节）。

\begin{Lexample}{vertpos}{\feat{VerticalPosition} 特性。}
  \fontspec{LibreCaslonText-Regular.otf}[VerticalPosition=Superior]
   Superior: 1234567890                                   \\
  \fontspec{LibreCaslonText-Regular.otf}[VerticalPosition=Numerator]
   Numerator: 12345                                       \\
  \fontspec{LibreCaslonText-Regular.otf}[VerticalPosition=Denominator]
   Denominator: 12345                                     \\
  \fontspec{LibreCaslonText-Regular.otf}[VerticalPosition=ScientificInferior]
   Scientific Inferior: 12345
\end{Lexample}

\pkg{realscripts} 宏包（在 \XeTeX\ 下也由 \pkg{xltxtra} 加载）
重定义了 \cmd\textsubscript\ 与 \cmd\textsuperscript\ 命令，
使其自动使用上述字体特性，包括在脚注标签中使用。若这是你想使用的
\pkg{xltxtra} 唯一功能，可考虑单独加载 \pkg{realscripts} 宏包。



\subsection{中日韩特性}

本节汇总主要面向中文、韩文与日文排版的特性。


\subsubsection{圈码 --- \texttt{nalt} }

有些字体配备了各种形式的丰富数字与数词。它们通过 \feat{Annotation}
特性（OpenType 特性 |nalt|）访问，按数值选取。注意编号从 0 开始。
\begin{Verbatim}
  \fontspec{Hiragino Maru Gothic Pro}
   1 2 3 4 5 6 7 8 9
  \def\x#1{\\{\addfontfeature{Annotation=#1}
            1 2 3 4 5 6 7 8 9 }}
  \x0\x1\x2\x3\x4\x5\x6\x7\x7\x8\x9
\end{Verbatim}


\subsubsection{字符宽度}\label{sec:CharacterWidth}

许多亚洲字体配备了不同间距的字符，以便挤进通常等宽的文字中。
这些通过 \feat{CharacterWidth} 特性访问。

\begin{features}{CharacterWidth}
\otf*{Proportional}{pwid}
\otf*{Full}        {fwid}
\otf*{Half}        {hwid}
\otf*{Third}       {twid}
\otf*{Quarter}     {qwid}
\otf*{AlternateProportional}{palt}
\otf*{AlternateHalf}{halt}
\cmidrule{2-4}
\otf{ResetAll}{}
\end{features}

日文假名（平假名或片假名）字形既可按比例排版，以更好地适应横向长度；
也可等宽排版，以纳入表意文字排版所要求的刚性网格。后一种情形下还有
半宽形式，用于在给定空间块内容纳更多假名（假名比周围的汉字简单）。
日文字体中的罗马字母也提供同样的特性，以便按与周围文字相同的风格
排版外来词。在找到支持这些特性的开源字体之前，示例暂缺。
%\begin{Verbatim}
%  \def\texta{ようこそ}\def\textb{ワカヨタレソ}
%  \def\test{\makebox[2cm][l]{\texta}%
%            \makebox[2.5cm][l]{\textb}%
%            \makebox[2.5cm][l]{abcdef}}
%  \fontspec{Hiragino Mincho Pro}
%  {\addfontfeature{CharacterWidth=Proportional}\test}\\
%  {\addfontfeature{CharacterWidth=Full}\test}\\
%  {\addfontfeature{CharacterWidth=Half}\test}
%\end{Verbatim}




\subsubsection{中日韩字形}

\begin{features}{CJKShape}
\otf{Traditional}{trad}
\otf{Simplified} {smpl}
\otf{JIS1978}    {jp78}
\otf{JIS1983}    {jp83}
\otf{JIS1990}    {jp90}
\otf{Expert}     {expt}
\otf{NLC}        {nlck}
\end{features}

关于中日韩表意文字“应当”长什么样，历来有多种标准。一些字体包含大量替代字形，
以便在合适的任一种形式下正确显示这些字符。\AAT\ 与 OpenType 字体都支持下列
\feat{CJKShape} 选项：\opt{Traditional}、\opt{Simplified}、\opt{JIS1978}、
\opt{JIS1983}、\opt{JIS1990} 与 \opt{Expert}。OpenType 还支持 \opt{NLC} 选项。

\begin{Xexample}[firstline=2]{ot-cjk-shape}{中日韩表意文字呈现的不同标准。}
  \LARGE\def\text{\noindent 唖噛躯 妍并訝}
  \fontspec{HaranoAjiGothic-Regular.otf}
  {\addfontfeature{CJKShape=Traditional}
  \text }                          \\
  {\addfontfeature{CJKShape=NLC}
  \text }                          
\end{Xexample}


\subsubsection{竖排}

\begin{features}{Vertical}
\otf*{RotatedGlyphs}         {vrt2}
\otf*{AlternatesForRotation} {vrtr}
\otf*{Alternates}            {vert}
\otf*{KanaAlternates}        {vkna}
\otf*{Kerning}               {vkrn}
\otf*{AlternateMetrics}      {valt}
\otf*{HalfMetrics}           {vhal}
\otf*{ProportionalMetrics}   {vpal}
\cmidrule{2-4}
\otf{ResetAll}{}
\end{features}

OpenType 提供了大量特性，以容纳竖排（中日韩及其他文字）所需的各种情形。
但 \pkg{fontspec} 尚未提供实现此类竖排的功能；如有可改进之处，欢迎联系。


\end{document}

% /©
% ------------------------------------------------
% The FONTSPEC package  <latex3.github.io/fontspec>
% ------------------------------------------------
% Copyright  2022-2026  The LaTeX project,  LPPL "maintainer"
% Copyright  2004-2022  Will Robertson
% Copyright  2009-2015  Khaled Hosny
% Copyright  2013       Philipp Gesang
% Copyright  2013-2016  Joseph Wright
% ------------------------------------------------
% This package is free software and may be redistributed and/or modified under
% the conditions of the LaTeX Project Public License, version 1.3c or higher
% (your choice): <http://www.latex-project.org/lppl/>.
% ------------------------------------------------
% ©/
